\chapter{Introduccion}
\label{cap:introduction}

\section{Contexto: Fatiga Mental y Fisica}

La fatiga es un estado fisiologico y psicologico complejo que afecta significativamente
el rendimiento cognitivo, la capacidad fisica y la calidad de vida. En contextos
academicos, laborales y deportivos, entender y medir la fatiga es crucial para
optimizar el desempenho y prevenir problemas de salud.

La fatiga se puede clasificar en dos categorias principales:

\begin{itemize}
    \item \textbf{Fatiga Fisica}: Cansancio o agotamiento muscular resultante del
          esfuerzo fisico intenso o prolongado. Se asocia con cambios en el metabolismo,
          acumulacion de lactato y deplecion de glucogeno.

    \item \textbf{Fatiga Mental}: Cansancio cognitivo resultante de la concentracion
          prolongada o tareas cognitivamente demandantes. Se asocia con cambios en
          actividad cerebral, control ejecutivo y vigilancia.
\end{itemize}

\section{Importancia del Monitoreo Multimodal}

Durante decadas, los investigadores han intentado entender la fatiga mediante
distintos enfoques:

\begin{enumerate}
    \item \textbf{Mediciones Subjetivas}: Cuestionarios y escalas de autopercepcion
          (ej: Escala Visual Analogica de Fatiga). Ventaja: economicas. Desventaja:
          susceptibles a sesgo y variabilidad individual.

    \item \textbf{Mediciones Fisiologicas Aisladas}: Frecuencia cardiaca,
          electroencefalograma (EEG), o respuesta electrodermal medidas
          independientemente. Ventaja: objetivas. Desventaja: incompletas.

    \item \textbf{Monitoreo Multimodal Simultaneo}: Multiples sensores wearables
          registrando simultaneamente diferentes variables fisiologicas, comportamentales
          y cognitivas. Ventaja: comprensiva, registra interacciones entre sistemas.
          Desventaja: mayor complejidad de analisis.
\end{enumerate}

El enfoque multimodal es particularmente importante porque la fatiga no es un
fenomeno unidimensional. Por ejemplo:

\begin{itemize}
    \item La frecuencia cardiaca puede aumentar o disminuir segun el tipo y
          intensidad de actividad y el estado cardiovascular del individuo.
    \item La actividad electrodermal refleja arousale (activacion del sistema nervioso
          simpatico) pero puede estar influenciada por factores emocionales o
          termicos.
    \item La actividad cerebral (EEG) proporciona informacion sobre estados cognitivos
          (alerta vs. somnolencia) pero requiere interpretacion cuidadosa de bandas
          de frecuencia.
\end{itemize}

Al monitorear multiples modalidades simultaneamente, podemos:
\begin{itemize}
    \item Detectar \textbf{patrones complejos} que emergen de la interaccion entre sistemas
    \item \textbf{Validar hallazgos} comparando senales independientes
    \item Desarrollar \textbf{modelos predictivos} mas robustos y generalizables
    \item Identificar \textbf{diferencias individuales} en como cada persona manifiesta la fatiga
\end{itemize}

\section{FatigueSet: Su Importancia en la Investigacion}

Aunque existen datasets de EEG, de actividad cardiaca, o de tareas cognitivas de
forma aislada, los datasets verdaderamente multimodales y accesibles publicamente
son escasos. FatigueSet llena este vacio proporcionando:

\begin{itemize}
    \item Una \textbf{recopilacion completa} de 4 dispositivos wearables de marcas
          comerciales diferentes pero bien documentadas
    \item Un \textbf{protocolo experimental riguroso} con contrabalanceo para
          investigar fatiga a multiples intensidades
    \item \textbf{Tareas cognitivas estandarizadas} (CRT, N-Back, Task Switching)
          junto con autoreportes validados
    \item \textbf{Datos procesados y limpios} listos para analisis, junto con codigo
          y documentacion
\end{itemize}

Este dataset ha permitido y seguira permitiendo investigacion en:
\begin{itemize}
    \item Deteccion automatica de fatiga mediante machine learning
    \item Entendimiento de mecanismos fisiologicos subyacentes a fatiga mental vs fisica
    \item Individualizacion de perfiles de respuesta a fatiga
    \item Validacion de nuevos sensores wearables
    \item Estudios de recuperacion y resilencia ante fatiga
\end{itemize}

\section{Objetivos de Este Documento}

Este documento tiene como objetivo proporcionar una referencia completa y autocontenida
del dataset FatigueSet. No asume conocimiento previo del dataset pero mantiene
rigor academico y tecnico.

Despues de leer este documento, el lector podra:

\begin{enumerate}
    \item Comprender la estructura general del estudio que genero el dataset
    \item Identificar que dispositivo recopilo que tipo de senales
    \item Interpretar cada una de las 19 variables fisiologicas disponibles
    \item Entender los patrones de fatiga observados en los datos
    \item Cargar y utilizar el dataset en sus propios analisis
    \item Reproducir los analisis mostrados en este documento
\end{enumerate}

\section{Estructura de Este Documento}

El documento esta organizado en 8 capitulos progresivos:

\begin{description}
    \item[Capitulo 1] (Este capitulo) - Motivacion y contexto
    \item[Capitulo 2] - Vision general del dataset y protocolo experimental
    \item[Capitulo 3] - Los 4 dispositivos wearables utilizados
    \item[Capitulo 4] - Las 5 tareas cognitivas y metricas conductuales
    \item[Capitulo 5] - Las 19 variables fisiologicas extraidas (enfoque en patrones de fatiga)
    \item[Capitulo 6] - Analisis cuantitativos y 14 visualizaciones de datos
    \item[Capitulo 7] - Interpretacion de hallazgos y signos de fatiga
    \item[Capitulo 8] - Guia practica para utilizar el dataset
\end{description}

Los apendices incluyen especificaciones tecnicas completas, instrumentos de encuesta,
codigo de ejemplo y un glosario de terminos fisiologicos y tecnicos.
